\documentclass{article}
\usepackage[utf8]{inputenc}
\usepackage[spanish]{babel}
\usepackage{graphicx}
\usepackage{verbatim}
\usepackage{moreverb}
\usepackage{amsmath}
\usepackage{amsfonts}
\usepackage{amssymb}
\usepackage{fancybox}
\usepackage{float}
\usepackage{fancyvrb}
\usepackage{color}
\usepackage{hyperref}
\usepackage{multirow}


\usepackage{anysize}
\marginsize{1.5cm}{1.5cm}{1cm}{1cm}

\renewcommand{\shorthandsspanish}{}

\newcommand{\HRule}{\rule{\linewidth}{0.5mm}}

\begin{document}

\thispagestyle{empty}

%%%%%%%%%%%%%%%%%%%%%%%%%%%%%%%%% PORTADA %%%%%%%%%%%%%%%%%%%%%%%%%%%%%%%%%%%%%%

\begin{titlepage}
\begin{center}

%Espacio antes del logo del itba
\Large \  \\[1.5cm]

\includegraphics[scale=0.40]{Imagenes/logo_itba}\\[1cm]
\textsc{\LARGE Criptografía y Seguridad}\\[1.5cm]
\textsc{\Large Trabajo práctico $\text{N}^{\circ}$2}\\[0.5cm]

\HRule \\[0.4cm]
{ \huge \bfseries Secreto Compartido en Imágenes con Esteganografía}\\[0.4cm]
\HRule \\[1.5cm]

\Large Autores: \\ [0.25cm]
\begin{tabular}{l @{\ \ -\ \ }l}
\Large Pablo Ballesty & \Large 49359\\[0.2cm]
\Large Nicolás Magni & \Large 48008\\[0.2cm]
\end{tabular}



\vspace{1cm}

\vfill
% La fecha queda abajo.
{\large \today}

\end{center}
\end{titlepage}

%%%%%%%%%%%%%%%%%%%%%%%%%%%%%%%%%%%%%%%%%%%%%%%%%%%%%%%%%%%%%%%%%%%%%%%%%%%%%%%%

\abstract{

El objetivo del presente informe es detallar las deciones tomadas durante el diseño y la implementación del Trabajo Práctico.
Así también, se presentan los resultados obtenidos a partir de las imágenes que envió la cátedra y se realiza un análisis de
los mismos. Se detallan las dificultades encontradas, y las soluciones propuestas. También se presentan conclusiones y posibles
extensiones del algoritmo.

}

\section{Introducción}

El Trabajo Práctico consiste en la implementación de un algoritmo de secreto compartido, aplicado a imágenes. Los algoritmos
de secreto compartido, consisten en ocultar información sensible, la cual llamaremos \emph{el secreto}, utilizando $n$ claves, de las cuales, para la
recuperación del secreto, serán necesarias $k$ (con $k \le n$).\\

Como en este caso, las claves son imágenes, al momento de generar las mismas, se utiliza una técnica de \emph{esteganografía} con el fin
de que el ojo humano no note, que se trata de claves.\\

El algoritmo que se utiliza en el trabajo se extrajo de la publicación \textbf{\emph{``Improvements in Geometry-Based Secret Image Sharing Approach
with Steganography``}} realizada por \textbf{Mustafá Ulutas}, \textbf{Vasif V. Nabiyev} y \textbf{Guzin Ulutas}, de la Universidad de Karadeniz, Turquía.

\section{Desarrollo}

La implementación del algoritmo no resulta difícil de implementar, sin embargo, precisa de cuidado en ciertos aspectos, para obtener resultados
exitosos. A continuación, se detallan las dificultades que presenta el algoritmo, y como fueron resultas en la práctica.

\subsection{Sistemas compatibles determinados}

\subsection{Validación de la distribución}


\end{document}